\chapter{LITERATURE REVIEW}
In this chapter, we summarize the works which are related to the case study of the elliptic curve cryptography.\\ 
A lot of research has been conducted in this field and we intern to review and analyze the information or data gathered so far.\\\\
\cite{Narayan2018} presented the order of the related algebraic group that was significantly easier to determine than in the general case for a family of elliptic curves for cryptographic purposes. This made it easier to meet the cryptographic requirement of a big prime factor for this arrangement. Another merit of this family was that group operations are slightly simplified.\\\\
\cite{Narayan2018a} presented a detailed overview of cryptography method that was based on complex mathematical problems such as prime number factorization and the discrete logarithm problem for elliptic curves. The method of continuously adding a point to an elliptic curve was known as elliptic curve point multiplication. It is used to create a trapdoor function in elliptic curve cryptography (ECC).\\\\
\cite{Koblitz2000} looked over the growth of elliptic curve cryptosystems. In 1976 Diffie and Hellman introduced the public-key cryptography. It was found out that the discrete logarithm problem had the potential to be used in public-key cryptosystems. Diffie and Hellman described the discrete logarithm problem as the challenge of finding logarithms with regard to a generator in the multiplicative group of the integers modulo $a$ prime, the discrete logarithm problem had since been redefined. This idea could be widened to arbitrary groups and, in particular, to elliptic curve groups. The resulting public-key systems had a small block size, are fast, and secure.\\\\
\cite{Koblitz1987} discussed public key cryptosystems that used the multiplicative group of a finite field had elliptic curves over finite fields. This is because the elliptic curve equivalent of the discrete logarithm problem was anticipated to be harder than the conventional discrete logarithm problem, especially over $GF(2^{n})$, these elliptic curve cryptosystems may be more safe on the nonsmoothness of the order of the cyclic subgroup created by a global point.\\\\
\cite{Smart1999} proposed a basic technique that addressed the discrete logarithm issue on trace one elliptic curves, which led to a linear algorithm. In practice, this meant that when elliptic curves were selected for cryptography, all curves with group orders equal to the finite field's order must be eliminated.\\\\
\cite{Barbulescu2015} investigated a field $GF(p_{n})$, where $n$ was a significant integer bigger that $1$, the hardness of the discrete logarithm issue was examined. The difficulty of this challenge was at the center of security evaluations for torus-based and pairing-based encryption, while being less studied than the tiny characteristic case or the prime field situation. The number field sieve(NFS) was the most well-known approach for tackling this problem. The ability to identify appropriate polynomials that specified the extension fields used in NFS was a significant component of this technique. Two new techniques for this problem was provided, each of which modifies the asymptotic complexity and allows for record-breaking calculations. It demonstrated these findings by computing discrete logarithms over a field $GF(p_{2})$ with cardinality of $180 digits (595 bits)$.
\cite{Chen2008} observed how the Legendre symbol behaved when applied to rational points on an elliptic curve, and came up with a family of binary sequences that had string pseudorandom features. That was, both the well-distribution measure and the correlation measure of order $k$ of the generated binary sequences are ``small" as measured by estimate of particular character sums along elliptic curves. It also offered a lower constraint on the linear complexity profile of these sequences. The findings suggested that such sequences behave similarly to Legendre sequences.\\\\
\cite{Koblitz2000} reviewed the growth on elliptic curve cryptosystems. Elliptic curve cryptography has seen increasing success in security applications in recent years due to features such as a small key size with a high level of protection. Their popularity has resulted in a wide range of implementations in terms of algorithms, curves, coordinate systems, platforms, and so forth.\\\\
\cite{Vidhya2019} highlighted works that overcame the drawback of the GCD attack in RSA. One of the most important varieties of Asymmetric cryptography was Rivest, Shamir, and Adleman (RSA). When compared to other cryptography techniques, the ky size in RSA is larger. Any two prime numbers were used to create the keys, with the Greatest Common Divisor (GCD) attacks, an unauthorized user might simply obtain the keys . One of RSA cryptography's major flaws was with a small key size. Elliptic curve cryptography (ECC) is also asymmetric cryptography. The goal of the proposed effort is to solve the GCD attack RSA flaw. Employing an exclusive OR(XOR), the RSA private and public keys were coupled with the ECC private and public keys, and new hybrid private and public keys were created. A new hybrid private key could be used to decrypt RSA data. The suggested method has the advantage of creating three tiers of cryptographic keys, which can avoid the GCD attack in RSA. When compared to the RSA algorithm's straight forward key creation, the key strength is greater.\\\\
\cite{Sutikno1998} discussed  the implementation of elliptic curves cryptosystem. After Diffie and Hellman introduced public key cryptography, T. ElGamal introduced a discrete logarithm problem-based public key cryptosystem approach in 1985. Neil Koblitz and Victor Miller proposed elliptic curve cryptosystems for the first time in 1985, on their own. This is because they used elliptic curve groups for arithmetic, elliptic curve cryptosystems are unique. The discrete logarithm issue in a set of points on an elliptic curve formed over a finite field is the basis for this cryptosystem. In an elliptic curve group, the discrete logarithm issue looked to be significantly more difficult than in other groups.\\\\
\cite{Schroeppel1995} explained how to build the Diffie-Hellman key exchange algorithm using a set of points on an elliptic curve over the field $F_{2n}$. A software version of this that used $n = 155$ might be optimized to achieve slightly quicker computation rates with a similar level of security than non-elliptic curve versions. The highly eddicient implementation was made possible by $F_{2n}$ quick calculation of reciprocals.\\\\
\cite{Kumar2016} proposed a novel image security algorithm based on elliptic curve cryptography (ECC) and DNA encoding, with the incressed usage of media in communications, image encryption was required to protect against attacks. The RGB image was first encoded using DNA encoding, then asymmetric encryption using elliptic curve Diffie-Hellman encryption  was applied (ECDHE). For analysis, the suggested technique was used to standard test photos. Key spacing, key sensitivity, and statistical analysis were all examined. According to the findings, the suggested method can withstand exhaustive attacks and is suitable for real applications.\\\\
\cite{Bellare2000} addressed the security of public-key cryptosystems in a ``multi-user'' scenario, specifically in the presence of assaults utilizing the encryption of related messages using distinct public keys, as demonstrated by Hastad`s classical RSA attacks. They demonstrated that security in a single -user context entails security in a multi-user setting if the former is regarded in the broad sense of ``indistinguishability'', pinning down a number of methods that are guaranteed to be secure against hastad-type assaults.\\\\
\cite{ElGamal1985} offered a novel signature technique, as well as a Diffie-Hellman key distribution scheme implementation that resulted in a public key cryptosystem. The difficulty of computing discrete logarithms over finite fields provided the security for both systems.\\\\
 \cite{mahto2016security}  described data transmission technologies that were transformed, as a result of this. It had provided platform for communicative parties to swiftly communicate information while simultaneously providing an opportunity for adversaries to strike. Cryptographic techniques were used to guarantee secrecy, integrity, and authentication to insecure data at rest, compares and contrasts the security of two widely used public-key cryptography systems, rivest shamir adleman (RSA) and elliptic curve cryptography (ECC). Rivest shamir adleman (RSA) was the first-generation public-key cryptography that was immensely popular since its introduction, while ECC has only lately gained prominence. The integer factorization issue (IFP) was used to secure the RSA cryptosystem, while the elliptic curve discrete logarithm problem (ECDLP) was used to secure the ECC cryptosystem. The major merit of ECC over RSA was because the best-known technique for solving the  elliptic curve discrete logarithm problem (ECDLP) requires full exponential time, whereas RSA's (IFP) takes sub-exponential time. This means that ECC can employ far fewer parameters than RSA while maintaining the same level of security. For instance the RSA technique requires a key to produce 112 bits of protection. ECC requires a key size of 224-255 bits, while RSA requires 2048 bits.\\\\
\cite{Danger2013} summarized elliptic curve encryption is subject to side-channel attacks. These attacks took use of biases in a variety of leakages, including power consumption, electromagnetic emission, execution time, and so on. To counter known assaults, countermeasures must be integrated. Because no one countermeasure can protect against all types of attacks, many of them must be integrated. However, because each has a hefty price tag, they can not all be used at the same time. Depending on the situation and the trade-off between security and performance, it is crucial to choose countermeasures carefully. The side-channel attacks on elliptic curve cryptography are summarized in this work, along with countermeasures. The cost in terms of time and space is stated for each countermeasure. Some assaults, such as the doubling attack, are clarified; others, such as the horizontal SVA, are enhanced; and new attacks, such as the horizontal attack against unified formulae, are detailed.